\section{optimisation}

We propose here an inexact CG-newton method for solving the optimisation problem: %using the ALM for constraint handling. 
\begin{equation}
\min_{p} \hat J(p),
\end{equation}
where
\begin{equation}
    \hat J(p) = J(u(p),p).
\end{equation}
Consider first the taylor series expansion of the objective \( \hat J \) around a point \( p_k \):
\begin{equation}
\hat J(p_k + \tilde{p}_k)
= \hat J(p_k)
+ \frac{d\hat J}{dp}(p_k)^{\top} \tilde{p}_k
+ \frac{1}{2} \tilde{p}_k^{\top} H(p_k) \tilde{p}_k
+ \mathcal{O}(\vert \tilde{p}_k \vert^{3}).
\end{equation}
Ignoring the $\mathcal{O}(\vert \tilde{p}_k \vert^{3})$ terms, we find $\tilde{p}_k$ that minimises this quadratic expression by setting its gradient with respect to $\tilde{p}_k$ to zero to obtain the Newton equation:
\begin{equation}
H(p_k)\tilde{p}_k = - \frac{d\hat J}{dp}(p_k).
\end{equation}
To avoid forming the full hessian matrix, we use the CG method to iteratively solve this system, which instead only requires the hessian-vector product \( H(p_k) v \) for some vector \( v \). 
We want to choose an appropriate forcing sequence $\eta_k$ to terminate the CG iterations when the residual \( r_k = H(p_k)\tilde{p}_k + \frac{d\hat J}{dp}(p_k) \) satisfies:
\begin{equation}
\| r_k \|_2 \leq \eta_k \left\| \frac{d\hat J}{dp}(p_k) \right\|_2.
\end{equation}
We use the choice: 
\begin{equation}
\eta_k = \min \left( 0.5, \sqrt{\left\Vert \frac{d\hat J}{dp}(p_k) \right\Vert} \right).
\end{equation}
We can prove (see e.g. \cite{nocedal2006numerical}) that choosing $\eta_k = O\left(\left\Vert \frac{d\hat J}{dp}(p_k) \right\Vert\right)$ leads to linear convergence for the iterations:
\begin{equation}
    \left\Vert \frac{d\hat J}{dp}(p_{k+1}) \right\Vert \leq C \left\Vert \frac{d\hat J}{dp}(p_k) \right\Vert,
\end{equation}
and that choosing $\eta_k = O\left(\left\Vert \frac{d\hat J}{dp}(p_k) \right\Vert^2\right)$ leads to quadratic (local) convergence:
\begin{equation}
    \left\Vert \frac{d\hat J}{dp}(p_{k+1}) \right\Vert \leq C \left\Vert \frac{d\hat J}{dp}(p_k) \right\Vert^2.
\end{equation}
albeit this tighter restriction may require too many CG iterations. The chosen $\eta_k$ is a compromise between these two extremes. 
Since the matrix is not necessarily positive definite, we also terminate if a direction of negative curvature is found \cite{nocedal2006numerical}.
\begin{equation}
\tilde{p}_k^T H(p_k)\tilde{p}_k < 0.
\end{equation}


